\section{Internalization in intuitionistic primitive recursive arithmetic}

\subsection{The arithmetic data types}

Fix once and for all a primitive-recursive coding of finite tuples and finite
lists by natural numbers.  Integers are represented by signed pairs and a
rational by a pair $(a,b)$ with $b>0$, modulo the decidable relation
$ad=cb$.  We do not need to form a quotient inside arithmetic: every theorem
is phrased in terms of representatives, and invariance under replacement by
an equivalent representative is proved by cross multiplication.  Addition,
multiplication, negation, conditional inversion, tuple projection, list
folds, and coefficient evaluation are primitive-recursive operations.

The following predicates are therefore quantifier-free after the defining
primitive-recursive functions have been named:
\begin{align*}
 \Rat(r)&: \text{$r$ is a rational code with positive denominator},\\
 \operatorname{Pt}(P)&: \text{$P$ is a nonzero rational triple},\\
 \operatorname{Ln}(\ell)&: \text{$\ell$ is a nonzero rational triple},\\
 \Inc(P,\ell)&: \ell\mathbin{\cdot}P=0,\\
 \Col(P,Q,R)&: \det(P,Q,R)=0,\\
 \Con(\ell,m,n)&: \det(\ell,m,n)=0.
\end{align*}
Apartness $x\mathrel{\#}y$ is the decidable assertion that the
cross-multiplied numerators differ.  This is the positive nondegeneracy notion
used before cancellation.  For rational codes, $x\#0$ computes a numerator
that is nonzero; multiplication cancellation reduces to the corresponding
integer lemma, which is proved in primitive recursive arithmetic by
quantifier-free induction.

\subsection{Shape of the eight arithmetical theorems}

For each fixed-size theorem, collect all input coordinates into one code
$d$.  Its internal statement has the form
\begin{equation}\label{eq:internal-fixed-shape}
 \forall d\;\bigl(\operatorname{WF}(d)\wedge
                    \operatorname{ND}(d)\wedge
                    \operatorname{Hyp}(d)
             \to \operatorname{Conclusion}(d)\bigr).
\end{equation}
Here $\operatorname{WF}$ checks rational and tuple codes,
$\operatorname{ND}$ is the finite conjunction of the explicit apartness
hypotheses, and every remaining atom is an equality of primitive-recursive
integer terms.  When the conclusion is existential, the proof supplies a
primitive-recursive witness function $w(d)$ and verifies its graph.  Examples
are the Ceva point, the Carnot conic coefficients, the Pappus line, and the
Brianchon concurrency point.

Carnot's polygon theorem has variable finite input.  Its internal statement
also receives the lists of side parameters and the coefficient-list
certificates
\[
 g_i(T)=c_i\prod_j(T-t_{ij}).
\]
The equality here means equality of finite coefficient lists.  The proof is
two nested primitive recursions: one evaluates the root product on a side and
the other telescopes the endpoint values around the polygon.  The induction
formula is an equality of coded integers, so it lies exactly in the
quantifier-free induction fragment of $\ipra$.

\subsection{Why fixed polynomial certificates are available in \texorpdfstring{$\ipra$}{iPRA}}

A fixed multivariate polynomial identity over $\Z$ can be verified by a
finite sequence of associativity, commutativity, distributivity, and numeral
computations.  Such a sequence is itself a finite derivation.  Equivalently,
one may implement a sparse polynomial normalizer by primitive recursion on
syntax and prove, by quantifier-free induction on that syntax, that evaluation
is preserved.  Applying the normalizer to the finitely many identities in
this report gives derivations of:
\begin{itemize}
\item all cross-product incidence equations;
\item the Pappus harmonic coordinates and bracket equation;
\item the Ceva and Menelaus determinants;
\item the three Carnot conic factorizations;
\item the affine four-point identity used by Hoehn;
\item the Pappus linear dependence; and
\item Pascal's three-coordinate identity.
\end{itemize}
The accompanying Python checker is not part of the logical proof.  It is an
independent way to detect transcription errors in precisely those fixed
identities.  A native formalization replaces each successful check by the
normalizer's proof object.

\subsection{Constructive normalizations}

The normalizations used in the human proofs do not conceal choice.  An
inverse of a $2\times2$ or $3\times3$ rational matrix is computed from its
adjugate once a named determinant is apart from zero.  When one of finitely
many minors must be nonzero, bounded search tests them in a fixed order.
Decidable bounded search and case analysis are intuitionistically valid and
primitive recursive.

For Brianchon, a fully intrinsic implementation can proceed as follows.  Code
the nondegenerate conic by a symmetric matrix $S$ (after clearing the factor
$2$ in mixed terms).  A tangent covector $\ell$ has contact point
$\operatorname{adj}(S)\ell$, up to a nonzero scalar.  From one contact point
and its tangent, bounded minor search constructs a rational projective frame
in which the point is $(0:0:1)$ and the tangent is $x=0$.  The conic then has
form
\[
 xz+A x^2+Bxy+C y^2=0,
\]
with $C\#0$ by nondegeneracy.  Replacing $z$ by $z+Ax+By$ and rescaling $z$
transforms it to $xz-y^2=0$.  These are explicit rational matrix operations;
no square root is used.  If one of the six tangent parameters is the point at
infinity of the chosen chart, a bounded search through seven rational
M\"obius transformations chooses a chart avoiding all six.  Each point excludes
at most one candidate, so the finite search is certified by a seven-case
pigeonhole argument formalizable in $\ipra$.  A still cleaner Lean interface
may take the projective matrix and six homogeneous parameters as a supplied
certificate and prove separately that the intrinsic algorithm produces one.

\subsection{No hidden classical principles}

The proofs never infer a witness from a double negation.  Every construction
is a term, a bounded search over decidable tests, or input certificate data.
The distinction matters most for Carnot's polygon theorem: an arbitrary
rational polynomial need not come with a rational splitting list, and
$\ipra$ should not be asked to prove such a list exists.  Instead, the theorem
states exactly what classical ``finite intersections counted with
multiplicity'' means computationally: a list and a verifiable factorization.
Likewise, ratios are concluded only after their denominators have been shown
apart from zero; the underlying polynomial equality is proved first.

\subsection{Transfer through \texttt{ConciseIPRA}}

The current ProveIt development contains a Lean theorem
\code{pra\_iff\_teq} and its theorem-level specialization
\code{pra\_theorem\_iff}.  They express that the ordinary $\ipra$ theory
\code{PRAAx} and the scalar equational theory \code{TeqAx} are
bi-interpretable, with both term round trips provably equal to the original
term.  Consequently the recommended completion route is:
\begin{enumerate}
\item construct each coded geometry derivation in \code{PRAAx};
\item apply \code{pra\_theorem\_iff} to transport it to \code{TeqAx}; and
\item if desired, continue through the proved equivalences from
      \code{TeqAx} to the smallest native \texttt{ConciseIPRA} checker
      calculus.
\end{enumerate}
This avoids re-proving field algebra directly in every compressed native
syntax.  The mathematical content of the present report is therefore a
finite family of explicit $\ipra$ derivations plus one primitive-recursive
list induction; the remaining work is proof-object engineering, not a new
geometric argument.
